% Editorial front matter; reflects the mathematical symbols used in the current text.
% When a chapter is added, review this list against that chapter's actual notation.
\chapter*{List of Notation}
\addcontentsline{toc}{chapter}{List of Notation}
\markboth{List of Notation}{List of Notation}

The following symbols are used in the present lecture notes. An operator carries a hat when it is written in abstract form; its matrix in a chosen basis may be written without a hat. Context distinguishes scalar coordinates from the corresponding operators.

\begingroup
\footnotesize
\renewcommand{\arraystretch}{1.08}
\setlength{\LTpre}{6pt}
\setlength{\LTpost}{8pt}
\begin{longtable}{@{}>{\raggedright\arraybackslash}p{0.31\textwidth}>{\raggedright\arraybackslash}p{0.61\textwidth}@{}}
\toprule
\textbf{Symbol} & \textbf{Meaning}\\
\midrule
\endfirsthead
\toprule
\textbf{Symbol} & \textbf{Meaning (continued)}\\
\midrule
\endhead
\bottomrule
\endfoot
\multicolumn{2}{@{}l}{\textbf{States and Hilbert-space notation}}\\[2pt]
$\ket{a}$, $\bra{a}$ & Ket representing a state vector, and its Hermitian-adjoint bra.\\
$\ket{n}$, $\bra{n}$ & Basis ket and bra; indices label basis states.\\
$\braket{a}{b}$ & Inner product of two state vectors.\\
$\mel{a}{\op{F}}{b}$ & Matrix element of operator $\op{F}$.\\
$a_n=\braket{n}{a}$ & Component of $\ket{a}$ in the basis $\{\ket{n}\}$.\\
$\sum_n\ket{n}\bra{n}=\mathbf{1}$ & Completeness relation for a discrete orthonormal basis.\\
$\delta_{mn}$ & Kronecker delta: $1$ when $m=n$, otherwise $0$.\\
\addlinespace[4pt]
\multicolumn{2}{@{}l}{\textbf{Operators, observables, and transformations}}\\[2pt]
$\op{F}$, $\op{T}$ & Abstract operators; $F$ and $T$ also denote matrix representations.\\
$\op{F}^{\dagger}$, $F^{\dagger}$ & Hermitian adjoint of an operator or matrix.\\
$\lambda$, $f_n$ & An eigenvalue; $f_n$ is an eigenvalue of observable $\op{F}$.\\
$\langle F\rangle$ & Expectation value of the observable $\op{F}$.\\
$\Delta F$ & Fluctuation of a quantity about its expectation value, as defined in context.\\
$[A,B]=AB-BA$ & Commutator of two operators or matrices.\\
$\{A,B\}=AB+BA$ & Anticommutator.\\
$S_{im}=\braket{i}{m}$ & Matrix connecting the $I$ and $M$ bases in Chapter 1.\\
$S^{\dagger}=S^{-1}$ & Unitarity of the representation-transformation matrix $S$.\\
$\operatorname{tr}(T)$ & Trace of a matrix or operator where defined.\\
\addlinespace[4pt]
\multicolumn{2}{@{}l}{\textbf{Position, momentum, and physical constants}}\\[2pt]
$\ket{x}$, $\ket{p}$ & Position and momentum eigenkets in one dimension.\\
$\vect{x}$, $\vect{p}$ & Position and momentum vectors in three dimensions.\\
$\op{x}_i$, $\op{p}_j$ & Cartesian components of the position and momentum operators.\\
$\psi(x)=\braket{x}{\psi}$ & Wavefunction in the position representation.\\
$\delta(x-x')$, $\delta^{(3)}(\vect{x}-\vect{x}')$ & Dirac delta distributions in one and three dimensions.\\
$\hbar$ & Reduced Planck constant.\\
$\ii$ & Imaginary unit, $\ii^2=-1$.\\
$\dd x$, $\dd^3x$ & One- and three-dimensional integration elements.\\
$\op{T}(\dd\vect{x})$ & Operator effecting an infinitesimal spatial translation.\\
$\braket{\vect{x}}{\vect{p}}$ & Position--momentum overlap, with Fourier normalization $(2\pi\hbar)^{-3/2}$.\\
\addlinespace[4pt]
\multicolumn{2}{@{}l}{\textbf{Spin and Stern--Gerlach notation}}\\[2pt]
$s_x$, $s_y$, $s_z$ & Components of spin angular momentum.\\
$\vect{J}$, $\vect{M}$, $\vect{B}$ & Angular momentum, magnetic moment, and magnetic field.\\
$\vect{e}_i$ & Unit vector along a coordinate axis.\\
\end{longtable}
\endgroup